\section{Métriques de comparaison}
\label{chap:métriques}
De nombreuses métriques existent afin d'évaluer les modèles de ML, 
dans ce chapitre, nous détaillons les métriques utilisées afin de 
comparer et analyser les résultats obtenus.

\noindent Nous débuterons par les métriques évaluant les modèles de Poisson.
\subsection{Erreur Quadratique Moyenne}
L'Erreur Quadratique Moyenne (EQM, ou \textit{Mean Squared Error} : MSE en anglais) 
($MSE = \frac{1}{n} \sum(y - \hat{y})^2$) 
mesure la moyenne des carrés des écarts entre les valeurs prédites et les valeurs 
réelles. En élevant au carré les erreurs de prédiction, la MSE pénalise les 
erreurs importantes. 


\subsection{Erreur Absolue Moyenne (MAE) }

L'Erreur Absolue Moyenne ($MAE= \frac{1}{n} \sum|y - \hat{y}|$ ) est une métrique 
qui mesure, en moyenne, la magnitude des erreurs entre les valeurs prédites par 
le modèle et les valeurs 
réelles observées, sans tenir compte de la direction de l'erreur. 
La MAE est appréciée pour sa simplicité et son interprétation intuitive. Elle donne 
une estimation directe de l'erreur moyenne sur l'ensemble des prédictions. 

\subsection{Mean\_poisson\_deviance}
Elle mesure l'erreur entre les valeurs réelles $y$ et les prédictions 
$\hat{y}$ pour des données de comptage (0,1,2,\dots). Plus la valeur est petite, 
meilleur est le modèle.
Pour une observation :
\[
D(y, \hat{y}) = 2 \left( y \log\left(\frac{y}{\hat{y}}\right) - (y - \hat{y}) \right)
\]
Pour $n$ observations :
\[
\text{mean\_poisson\_deviance} = \frac{1}{n} \sum_{i=1}^{n} 2 \left( y_i \log\left(\frac{y_i}{\hat{y}_i}\right) - (y_i - \hat{y}_i) \right)
\]




\vspace*{1cm}
\noindent Dans la suite, $N$ désigne le nombre de matchs de l'échantillon, $o_i \in \{0, 1\}$ 
l'observation réelle et $p_i \in [0, 1]$ la probabilité prédite par le modèle.

\subsection{Score F1}
\label{sec:f1}
Le score F1 est un compromis entre la précision
$precision = \frac{TP}{TP + FP}$ et le rappel
$recall = \frac{TP}{TP + FN}$.

Il est défini par :
\[
F1 = 2 \times \frac{Precision \times Recall}{Precision + Recall}
\]

\noindent Il correspond à la moyenne harmonique de la précision et du rappel (Recall).  
Le rappel se concentre sur les faux négatifs, c'est-à-dire les vrais positifs que le modèles 
aurait raté lors de la prédiction. Tandis que la précision se concentre sur 
les faux positifs, c'est-à-dire la proportion d'erreur lors de la prédiction. 
Ainsi, le score F1 permet de combiner ces deux métriques en une, 
afin d'évaluer les performances globales du modèle, notamment dans le cas de 
classes déséquilibrées.
Dans notre projet, il est évalué au seuil de confiance optimal 
identifié pour chaque modèle, afin de refléter les meilleures performances retenues.
Le seuil de confiance optimal est trouvé en maximisant le profit basé sur le ROI.

\subsection{Return on Invest (ROI)}
\label{sec:roi}
Le ROI (\textit{Return on Investment}) est une métrique permettant d'évaluer la 
rentabilité d'un modèle de paris, défini par :
\[
ROI = \frac{gain - mise}{mise} \times 100
\]
Bien qu'il offre une mesure intuitive de la rentabilité, il ne tient pas compte 
du volume de paris effectués.

\subsection{Brier-Score}
\label{sec:BS}
Nous avons également utilisé le Brier Score afin de comparer les performances 
de nos modèles. Il s'agit d'une approche similaire à l'EQM, à la différence près 
qu'il s'applique aux probabilités prédites par un modèle de classification. Il mesure 
ainsi l'écart quadratique moyen entre la probabilité prédite et l'observation réelle, 
permettant d'évaluer la qualité des prédictions probabilistes du modèle. La formule 
la plus commune de ce score est :
\[
BS = \frac{1}{N} \sum_{i=1}^{N} (p_i - o_i)^2
\]
Plus ce score est faible, meilleure est la calibration du modèle.


\subsection{Log-loss}
\label{sec:log-loss}

La Log-loss est une fonction de coût évaluant la précision des probabilités prédites 
par nos classifieurs. Contrairement au Brier Score, elle pénalise plus sévèrement 
les prédictions confiantes mais erronées. Pour un problème binaire, 
elle est définie par :
\[
LogLoss = - \frac{1}{N} \sum_{i=1}^{N} \left[ o_i \ln(p_i) + (1 - o_i) \ln(1 - p_i) \right]
\]
Comme pour le Brier Score, un score faible indique un modèle performant.




